\documentclass[a4paper,14pt,oneside,final]{report}

\usepackage[utf8]{inputenc}
\usepackage[T2A]{fontenc}
\usepackage[english, russian]{babel}
\usepackage{multicol}
\usepackage{multirow}
\usepackage{setspace}
%\usepackage{pscyr}
%\renewcommand{\rmdefault}{cmr}
\usepackage[final,hidelinks]{hyperref}
\usepackage[authoryear,square,numbers,sort&compress]{natbib}
%\usepackage{cite}
%\usepackage{times}
\usefont{T2A}{ftm}{m}{sl}
\setlength{\bibsep}{0em}
\PassOptionsToPackage{hyphens}{url}

\usepackage{hyphenat}

\addto\extrasrussian{%
  \def\equationautorefname{формула}%
  \def\figureautorefname{рисунок}%
  \def\listingautorefname{листинг}%
  \def\tableautorefname{таблица}%
}

\addto\captionsrussian{
  \renewcommand\contentsname{\centerline{\bfseries\large{\MakeUppercase{содержание}}}}
  \renewcommand{\bibsection}{\sectioncentered*{Cписок использованных источников}}
  \renewcommand{\listingscaption}{Листинг}
}

\AtEndDocument{
  \addcontentsline{toc}{section}{Cписок использованных источников}%
}
\usepackage{extsizes}

\sloppy

\usepackage{microtype}
\newlength{\fivecharsapprox}
\setlength{\fivecharsapprox}{6ex}


\usepackage{indentfirst}
\setlength{\parindent}{\fivecharsapprox} 

% Зачем: pdflscape поворачивает страницу в результирующем PDF, чтобы широкие диаграммы читались без поворота головы
\usepackage{pdflscape}
\usepackage[left=3cm,top=2.0cm,right=1.5cm,bottom=2.7cm]{geometry}

\frenchspacing

\usepackage{perpage}
\MakePerPage{footnote}

\makeatletter 
\def\@makefnmark{\hbox{\@textsuperscript{\normalfont\@thefnmark)}}}
\makeatother

\usepackage[bottom]{footmisc}

\makeatletter
\renewcommand{\thesection}{\arabic{section}}
\makeatother

\setcounter{secnumdepth}{3}


% Зачем: Настраивает отступ между таблицей с содержанимем и словом СОДЕРЖАНИЕ
\usepackage{tocloft}
\setlength{\cftbeforetoctitleskip}{-1em}
\setlength{\cftaftertoctitleskip}{1em}


% Зачем: Определяет отступы слева для записей в таблице содержания.
\makeatletter
\renewcommand{\l@section}{\@dottedtocline{1}{0.5em}{1.2em}}
\renewcommand{\l@subsection}{\@dottedtocline{2}{1.7em}{2.0em}}
\renewcommand{\l@subsubsection}{\@dottedtocline{2}{3.7em}{2.5em}}
\makeatother


% Зачем: Работа с колонтитулами
\usepackage{fancyhdr} 
\pagestyle{fancy}


\fancyhf{} 
\fancyfoot[R]{\thepage}
\renewcommand{\footrulewidth}{0pt} 
\renewcommand{\headrulewidth}{0pt}
\fancypagestyle{plain}{ 
    \fancyhf{}
    \rfoot{\thepage}}


% Зачем: Задает стиль заголовков раздела жирным шрифтом, прописными буквами, без точки в конце
\makeatletter
\renewcommand\section{%
  \clearpage\@startsection {section}{1}%
    {\fivecharsapprox}%
    {-1em \@plus -1ex \@minus -.2ex}%
    {1em \@plus .2ex}%
    {\raggedright\hyphenpenalty=10000\normalfont\large\bfseries\MakeUppercase}}
\makeatother


% Зачем: Задает стиль заголовков подразделов

\makeatletter
\renewcommand\subsection{
  \@startsection{subsection}{2}
    {\fivecharsapprox}
    {-1em \@plus -1ex \@minus -.2ex}
    {1em \@plus .2ex}
    {\raggedright\hyphenpenalty=10000\normalfont\normalsize\bfseries}}
\makeatother


% Зачем: Задает стиль заголовков пунктов
\makeatletter
\renewcommand\subsubsection{
  \@startsection{subsubsection}{3}%
    {\fivecharsapprox}%
    {-1em \@plus -1ex \@minus -.2ex}%
    {1em \@plus .2ex}
    {\raggedright\hyphenpenalty=10000\normalfont\normalsize\bfseries}}
\makeatother

% Зачем: для оформления введения и заключения, они должны быть выровнены по центру.
\makeatletter
\newcommand\sectioncentered{%
  \clearpage\@startsection {section}{1}%
    {\z@}%
    {-1em \@plus -1ex \@minus -.2ex}%
    {1em \@plus .2ex}%
    {\centering\hyphenpenalty=10000\normalfont\large\bfseries\MakeUppercase}%
    }
\makeatother

\usepackage[final]{graphicx}
\DeclareGraphicsExtensions{.png,.jpg}

% Зачем: Добавление подписей к рисункам
\usepackage{caption}
\usepackage{subcaption}
\usepackage{needspace}

% Зачем: оформление фрагментов кода как рисунков без разрыва между листингом и подписью
% #1 -- резерв высоты, #2 -- опции listings, #3 -- имя файла в code-snippets/, #4 -- текст подписи, #5 -- метка
\newcommand{\trmcodefigure}[5][32]{%
    \Needspace{#1\baselineskip}%
    \begin{figure}[!ht]%
        \centering%
        \begin{minipage}{\linewidth}%
            \lstinputlisting[style=trmcode,#2]{code-snippets/#3}%
        \end{minipage}%
        \caption{Фрагмент кода #4}%
        \label{#5}%
    \end{figure}%
}

% Зачем: поворот ячеек таблиц на 90 градусов
\usepackage{rotating}
\DeclareRobustCommand{\povernut}[1]{\begin{sideways}{#1}\end{sideways}}

\DeclareRobustCommand{\x}[1]{\text{#1}}

% Зачем: Задание подписей, разделителя и нумерации частей рисунков
% сказали что надо без курсива
% \DeclareCaptionLabelFormat{stbfigure}{Рисунок \emph{#2}}
% \DeclareCaptionLabelFormat{stbtable}{Таблица \emph{#2}}
% \DeclareCaptionLabelFormat{stblisting}{Листинг \emph{#2}}
\DeclareCaptionLabelFormat{stbfigure}{Рисунок #2}
\DeclareCaptionLabelFormat{stbtable}{Таблица #2}
\DeclareCaptionLabelFormat{stblisting}{Листинг #2}
\DeclareCaptionLabelSeparator{stb}{~--~}
\captionsetup{labelsep=stb}
\captionsetup[figure]{labelformat=stbfigure, justification=centering, font={small}}
\captionsetup[listing]{labelformat=stblisting,justification=centering, font={small}}
% сказали что надо слева
% \captionsetup[table]{labelformat=stbtable, justification=centering, font={small}}
\captionsetup[table]{labelformat=stbtable, justification=raggedright,singlelinecheck=false, font={small}}
\renewcommand{\thesubfigure}{\asbuk{subfigure}}

% Зачем: Окружения для оформления формул
\usepackage{calc}
\newlength{\lengthWordWhere}
\settowidth{\lengthWordWhere}{где}
\newenvironment{explanation}
    {
    \begin{itemize}[leftmargin=0cm, itemindent=\lengthWordWhere + \labelsep , labelsep=\labelsep]

    \renewcommand\labelitemi{}
    }
    {
    \\[\parsep]
    \end{itemize}
    }


\usepackage{tabularx}

% Зачем: типы колонок для табличных окружений с автопереносом текста и выравниванием по \textwidth.
% Префикс \hspace{0pt} включает перенос слов начиная с первого слова ячейки.
\newcolumntype{L}[1]{>{\raggedright\arraybackslash\hspace{0pt}}p{#1}}
\newcolumntype{Y}{>{\raggedright\arraybackslash\hspace{0pt}}X}
\newcolumntype{C}{>{\centering\arraybackslash\hspace{0pt}}X}

\newenvironment{explanationx}
    {
    \noindent 
    \tabularx{\textwidth}{@{}ll@{ --- } X }
    }
    { 
    \\[\parsep]
    \endtabularx
    }

\usepackage{amsmath}

\usepackage{amsfonts}
\usepackage{amssymb}
\usepackage{amsthm}
\usepackage{calc}
\usepackage{fp}
\usepackage{enumitem}

\makeatletter
 \AddEnumerateCounter{\asbuk}{\@asbuk}{щ)}
\makeatother


\setlist{nolistsep}
\renewcommand{\labelenumi}{\asbuk{enumi})}
\renewcommand{\labelenumii}{\arabic{enumii})}


\setlist[itemize,0]{itemindent=\parindent + 2.2ex,leftmargin=0ex,label=--}
\setlist[enumerate,1]{itemindent=\parindent + 2.7ex,leftmargin=0ex}
\setlist[enumerate,2]{itemindent=\parindent + \parindent - 2.7ex}

% Зачем: Включение номера раздела в номер формулы. Нумерация формул внутри раздела.
\AtBeginDocument{\numberwithin{equation}{section}}

% Зачем: Включение номера раздела в номер таблицы. Нумерация таблиц внутри раздела.
\AtBeginDocument{\numberwithin{table}{section}}

% Зачем: Включение номера раздела в номер рисунка. Нумерация рисунков внутри раздела.
\AtBeginDocument{\numberwithin{figure}{section}}

% Зачем: Включение номера раздела в номер листинга. Нумерация листингов внутри раздела.
\AtBeginDocument{\numberwithin{listing}{section}}


\usepackage{makecell}
\usepackage{multirow}
\usepackage{array}


\usepackage{textcomp}

\usepackage{siunitx}
\sisetup{
  binary-units = true,
  output-decimal-marker = {,},
  per-mode = symbol,
  range-phrase = --,
}
\DeclareSIUnit{\sample}{S}


\newcommand{\ignore}[2]{\hspace{0in}#2}


\usepackage{verbatim}
\usepackage{xcolor}
\usepackage{minted}

% Зачем: оформление листингов кода с обрамлением, переносом длинных строк и подсветкой синтаксиса
\usepackage{listings}

\definecolor{trmCodeBackground}{RGB}{248,249,251}
\definecolor{trmCodeFrame}{RGB}{120,135,170}
\definecolor{trmCodeKeyword}{RGB}{20,75,170}
\definecolor{trmCodeComment}{RGB}{110,130,150}
\definecolor{trmCodeString}{RGB}{30,130,80}
\definecolor{trmCodeNumber}{RGB}{120,80,160}
\definecolor{trmCodeRule}{RGB}{120,130,150}

\lstdefinestyle{trmcode}{
    backgroundcolor=\color{trmCodeBackground},
    basicstyle=\ttfamily\footnotesize,
    keywordstyle=\color{trmCodeKeyword}\bfseries,
    commentstyle=\color{trmCodeComment}\itshape,
    stringstyle=\color{trmCodeString},
    numberstyle=\tiny\color{trmCodeRule},
    numbers=left,
    numbersep=8pt,
    frame=single,
    rulecolor=\color{trmCodeFrame},
    framerule=0.5pt,
    framesep=4pt,
    breaklines=true,
    breakatwhitespace=false,
    breakindent=12pt,
    postbreak=\mbox{\textcolor{trmCodeRule}{$\hookrightarrow$}\space},
    showstringspaces=false,
    tabsize=4,
    captionpos=b,
    abovecaptionskip=4pt,
    belowcaptionskip=4pt,
    xleftmargin=0pt,
    xrightmargin=0pt,
    aboveskip=10pt,
    belowskip=10pt,
    upquote=true,
    columns=fullflexible,
    keepspaces=true
}

\lstdefinelanguage{JavaScriptES}{
    keywords={await, async, break, case, catch, class, const, continue, default, delete, do, else, export, extends, false, finally, for, from, function, if, import, in, instanceof, let, new, null, of, return, static, super, switch, this, throw, true, try, typeof, undefined, var, void, while, with, yield},
    keywordstyle=\color{trmCodeKeyword}\bfseries,
    ndkeywords={React, useState, useEffect, useRef, useMemo, useCallback, dispatch, useSelector, useDispatch, useTranslation, console, fetch, Map, JSON, Math},
    ndkeywordstyle=\color{trmCodeKeyword},
    sensitive=true,
    comment=[l]{//},
    morecomment=[s]{/*}{*/},
    morestring=[b]',
    morestring=[b]",
    morestring=[b]`
}

\lstset{style=trmcode}

%\AtBeginDocument{\numberwithin{lstlisting}{section}}

\usepackage[normalem]{ulem}

\renewcommand{\UrlFont}{\small\rmfamily\tt}

% Магия для подсчета разнообразных объектов в документе
\usepackage{lastpage}
\usepackage{totcount}
\regtotcounter{section}

\usepackage{etoolbox}


\newcounter{totfigures}
\newcounter{tottables}
\newcounter{totreferences}
\newcounter{totequation}

\providecommand\totfig{} 
\providecommand\tottab{}
\providecommand\totref{}
\providecommand\toteq{}

\makeatletter
\AtEndDocument{%
  \addtocounter{totfigures}{\value{figure}}%
  \addtocounter{tottables}{\value{table}}%
  \addtocounter{totequation}{\value{equation}}
  \immediate\write\@mainaux{%
    \string\gdef\string\totfig{\number\value{totfigures}}%
    \string\gdef\string\tottab{\number\value{tottables}}%
    \string\gdef\string\totref{\number\value{totreferences}}%
    \string\gdef\string\toteq{\number\value{totequation}}%
  }%
}
\makeatother

\pretocmd{\section}{\addtocounter{totfigures}{\value{figure}}\setcounter{figure}{0}}{}{}
\pretocmd{\section}{\addtocounter{tottables}{\value{table}}\setcounter{table}{0}}{}{}
\pretocmd{\section}{\addtocounter{totequation}{\value{equation}}\setcounter{equation}{0}}{}{}
\pretocmd{\bibitem}{\addtocounter{totreferences}{1}}{}{}



% Для оформления таблиц не влязящих на 1 страницу
\usepackage{longtable}

\usepackage{gensymb}

% Зачем: преобразовывать текст в верхний регистр командой MakeTextUppercase
\usepackage{textcase}

%  Переносы в словах с тире \hyph.

\def\hyph{-\penalty0\hskip0pt\relax}

% Добавляем левый отступ для библиографии

\makeatletter
\renewenvironment{thebibliography}[1]
     {\sectioncentered*{Cписок использованных источников}
      \@mkboth{\MakeUppercase\refname}{\MakeUppercase\refname}%
      \list{\@biblabel{\@arabic\c@enumiv}}%
           {\settowidth\labelwidth{\@biblabel{#1}}%
            \setlength{\itemindent}{\dimexpr\labelwidth+\labelsep+1em}
            \leftmargin\z@
            \@openbib@code
            \usecounter{enumiv}%
            \let\p@enumiv\@empty
            \renewcommand\theenumiv{\@arabic\c@enumiv}}%
      \sloppy
      \clubpenalty4000
      \@clubpenalty \clubpenalty
      \widowpenalty4000%
      \sfcode`\.\@m}
     {\def\@noitemerr
       {\@latex@warning{Empty `thebibliography' environment}}%
      \endlist}
\makeatother

\newcommand{\intro}[3]{
    \stepcounter{section}
        \sectioncentered*{ПРИЛОЖЕНИЕ \MakeUppercase{#1}}
     \begin{center} 
        \bf{(#2)}\\
        \bf{#3}
    \end{center}
    \markboth{\MakeUppercase{#1}}{}
    \addcontentsline{toc}{section}{Приложение \MakeUppercase{#1} (#2) #3}
}

\phantomsection\pagebreak % исправляет нумерацию в документе и исправляет гиперссылки в pdf